\documentclass[10pt,a4paper]{amsart}
\usepackage[margin=19mm]{geometry}
\usepackage{amsmath,amssymb,mathtools}
\usepackage[scr=rsfs,cal=euler]{mathalfa}
\usepackage{xfrac}
\usepackage[unicode,hidelinks,bookmarks=false]{hyperref}
\newcommand{\ie}{i.e.\ }
\newcommand{\cf}{cf.\ }
\newcommand{\resp}{respectively\ }
\newcommand{\Subsection}{\subsection}
\providecommand{\nobreakdash}{\nobreak\hbox{-}}
\setlength{\emergencystretch}{2em}
\multlinegap0pt
\usepackage{mathtools}
\usepackage{amssymb}
\newtheorem{prop}{Proposition}
\newtheorem{cor}{Corollary}
\newtheorem{lem}{Lemma}
\newcommand{\HO}{\operatorname{H}}
\newcommand{\Nis}{\mathrm{Nis}}
\newcommand{\Zar}{\mathrm{Zar}}
\newcommand{\et}{\mathrm{\acute{e}t}}
\newcommand{\spec}{\operatorname{Spec}}
\title{Gersten-type conjecture for henselian local rings of normal crossing varieties}
\author{Makoto Sakagaito}
\date{Source: arXiv:2402.18042v4 (CC BY 4.0)}
\begin{document}
\maketitle

\noindent\emph{Source and licence.} Gersten-type conjecture for henselian local rings of normal crossing varieties. Version arXiv:2402.18042v4 (CC BY 4.0), distributed by arXiv under the Creative Commons Attribution 4.0 International licence, http://creativecommons.org/licenses/by/4.0/, as recorded in the arXiv licence metadata for this exact version and shown on the abstract page https://arxiv.org/abs/2402.18042v4. Full licence notice: \texttt{licenses/arxiv-2402.18042-CC-BY-4.0.txt}.\par\smallskip

\noindent\textit{Editorial scope.} Counted unit: the article's lem nrhenij (source0.tex lines 8696--8721). The article's complete original proof of it is delivered with it (source0.tex lines 8921--9015, one \texttt{proof} environment of 95 lines, which the article places away from the statement). No mathematics is written, repaired or re-derived: the statement and the proof are the article's own lines, in the article's own notation, and no retained line is substituted or re-flowed.  Retained uncounted as the source-local context the proof needs: lines 6435--6476; lines 8151--8167; lines 8518--8549; lines 8901--8911.  Nothing was omitted from any retained span, so the package carries no omission record.  No retained line contains a citation, so the wrapper carries no bibliography and no bibliography entry.  No retained line contains a figure environment, an image-inclusion command or drawing code, so no image is delivered and the package declares no asset dependency.  Editorial formatting only, in addition: an AMS \texttt{amsart} preamble that restates the article's own declarations and macros used by the retained lines, namely the environments \texttt{prop}, \texttt{cor}, \texttt{lem} on separate counters, together with the standard packages those lines need.  The retained environments print in this document as Proposition 1, Corollary 1, Lemma 1 in order of appearance, whereas the article prints the same items with the numbering of its own full document.  The complete native source of the article as distributed in the arXiv e-print for this exact version is bundled unchanged as \texttt{sources/155255/source0.tex}.
\begin{prop}\upshape\label{constm}
Let $A$ be a local ring of a normal crossing variety over the spectrum of  
a field $k$. Let $n\geq 0$, $l>0$ be integers with $(l, \operatorname{char}(k))=1$.
Then we have isomorphisms
%
\begin{equation}\label{Vanineg}
\mathcal{H}^{q}(\mathbb{Z}/l(n)_{\Zar}^{\spec(A)})=0    
\end{equation}
%
for $q<0$ and
%
\begin{align}\label{Zerods}
\HO^{0}_{\Zar}(
A, 
\mathbb{Z}/l(n)
)  
\xrightarrow{\simeq}
\bigoplus_{x\in\spec(A)^{(0)}}
\HO^{0}_{\Zar}(
\kappa(x),
\mathbb{Z}/l(n)
)
\simeq
\underbrace{
\mu_{l}^{\otimes n}
\oplus\cdots
\oplus
\mu_{l}^{\otimes n}
}_{\#\left((\spec(A))^{(0)}\right) ~ \textrm{times}}.
\end{align}
%
Moreover, we have an isomorphism
%
\begin{equation}\label{Vanipos}
\mathcal{H}^{q}(
\mathbb{Z}/l(n)_{\et}^{\spec(A)}
)    
=0
\end{equation}
%
for $q\geq\#\left((\spec(A))^{(0)}\right)$.
\end{prop}

\begin{align}\label{mfmtm}
&\tau_{\leq n+1}(
\mathbb{Z}/l(n)_{\et}^{\spec(C)}
)    
\simeq 
\mu_{l}^{\otimes n}
\to
i_{*}\mu_{l}^{\otimes n}  \nonumber 
\\
\to
&i_{*}\left(
\tau_{\leq 0}(\mathbb{Z}/l(n)_{\et}^{\spec(C/(\pi))}
)
\right)
\to 
i_{*}\mathbb{Z}/l(n)_{\et}^{\spec(C/(\pi))}
\end{align}

\begin{cor}\upshape\label{inclaffmt}
Let $B$ be a henselian discrete valuation ring of mixed characteristic $(0, p)$ and
$\pi$ a prime element of $B$. Let $n$ be a non-negative integer and 
$l$ an integer which is prime to $p$.
Let 
%
\begin{equation*}
C=
B[T_{0}, \cdots, T_{N}]/
(T_{0}\cdots T_{a}-\pi)
\end{equation*}
%
for 
$0\leq a\leq N$
and
$\alpha: \spec(C)_{\et}\to \spec(C)_{\Nis}$
the canonical map of sites. Then the homomorphism
%
\begin{equation*}
\Gamma\left(
\spec(C),
R^{n+1}\alpha_{*}\mathbb{Z}/l(n)_{\et}
\right)
\to
\Gamma\left(
\spec(C/(\pi)),
R^{n+1}\alpha_{*}\mathbb{Z}/l(n)_{\et}
\right)
\end{equation*}
%
is injective where the homomorphism is induced by (\ref{mfmtm}).
\end{cor}

\begin{equation}\label{probnc}
\Gamma\left(
\mathfrak{X},
R^{n+1}\alpha_{*}\mu_{l}^{\otimes n}
\right)
\to
\Gamma\left(
Y,
R^{n+1}\alpha_{*}\mu_{l}^{\otimes n}
\right)
\end{equation}

\begin{lem}\upshape\label{nrhenij}
Let $B$ be a discrete valuation ring and $\pi$ a prime element of $B$. Let $R$ 
be the henselization of the local ring of
%
\begin{math}
C=B[T_{0}, \cdots, T_{N}]
/(T_{0}\cdots T_{a}-\pi T_{N}^{b})    
\end{math}
%
at a point $x$ of $\spec(C)$
and $l$ an integer which is prime to $\operatorname{char}(B)$.
%
Then the homomorphism
%
\begin{equation}\label{homnr}
\HO^{n}_{\et}(R, \mu_{l}^{\otimes n})
\to
\displaystyle\bigoplus_{x\in \spec(R)^{(0)}}
\HO^{n}_{\et}(
\kappa(x),
\mu_{l}^{\otimes n}
)
\end{equation}
%
is injective for any integer $n\geq 0$.
\end{lem}

\begin{proof}\upshape
Let 
%
\begin{math}
j: D(T_{N+1})\to  
\mathfrak{X}
\end{math}
%
be an open immersion of $\mathfrak{X}$ such that
%
\begin{align*}
D(T_{N+1})
:=
\{
\mathfrak{p}\in \mathfrak{X}|
T_{N+1}\notin\mathfrak{p}
\}
=\spec(
B[S_{0}, \cdots, S_{N}]/
(S_{0}\cdots S_{a}-\pi)
).
\end{align*}
%
Then the homomorphism
%
\begin{equation*}
R^{n+1}\alpha_{*}\mu_{l}^{\otimes n}
\to 
j_{*}j^{*}R^{n+1}\alpha_{*}\mu_{l}^{\otimes n}
\end{equation*}
%
is injective by Lemma \ref{nrhenij}. So the homomorphism
%
\begin{align}\label{inclpta}
\Gamma\left(
\mathfrak{X},
R^{n+1}\alpha_{*}\mu_{l}^{\otimes n}
\right)
\to
\Gamma\left(
\spec(C),
R^{n+1}\alpha_{*}\mu_{l}^{\otimes n}
\right)
\end{align}
%
is injective where
%
\begin{equation*}
C=
B[S_{0}, \cdots, S_{N}]/
(S_{0}\cdots S_{a}-\pi).    
\end{equation*}
% 
Then the composite of the homomorphism (\ref{inclpta}) and the homomorphism
%
\begin{align*}
\Gamma\left(
\spec(C),
R^{n+1}\alpha_{*}\mathbb{Z}/l(n)
\right)
\to
\Gamma\left(
\spec(C/(\pi)),
R^{n+1}\alpha_{*}\mathbb{Z}/l(n)
\right)
\end{align*}
%
agrees with the composite of the homomorphism (\ref{probnc}) and
the composition
%
\footnotesize
\begin{align*}
\Gamma\left(
Y, 
R^{n+1}\alpha_{*}
\mu_{l}^{\otimes n}
\right) 
\to
\Gamma\left(
\spec(C/(\pi)),
R^{n+1}\alpha_{*}
\mu_{l}^{\otimes n}
\right)
\to
\Gamma\left(
\spec(C/(\pi)),
R^{n+1}\alpha_{*}
\mathbb{Z}/l(n)
\right)
\end{align*}
\normalsize
%
(cf. Proposition \ref{constm}).
Hence the statement follows from Corollary \ref{inclaffmt}.
\end{proof}

\end{document}
